\section{Hardness under the Unique Games Conjecture}
\label{sec:ugc}

We prove ordinary-UGC hardness below $4/3$ for treedepth and
treewidth, and below $8/7$ for total elimination depth. We begin
with the treedepth construction. Its first-branch lemmas will also
be used for total depth.

\subsection{Treedepth}
\label{subsec:ugc-treedepth}

For a fixed integer $q$, the construction replaces the source graph by
$q$ clique layers. A large independent set gives a shallow elimination forest: its copies
can be placed in separate branches after the remaining vertices have
been removed. For soundness, we prove that branching can save few
levels when the source independence number is small.

We use the following normalized source promise.

\begin{theorem}[Khot--Regev, normalized form {\cite[Theorem~3.2 and Section~4]{KhotRegev}}]\label{thm:kr}
For every fixed rational $0<\eta<1/10$, it is Unique Games-hard to
distinguish simple unweighted graphs $G$ on $n\geq1$ vertices with
$\alphaI(G)\geq(1/2-\eta)n$ \textup{(YES)} from those with
$\alphaI(G)<\eta n$ \textup{(NO)}.
\end{theorem}

Appendix~\ref{app:kr-normalization} derives this unweighted form from
the published weighted promise.

\subsubsection{The layered construction and first-branch notation}

Let $G$ be a finite simple graph on $n$ vertices and let $q\geq1$
be an integer. Define $S_q(G)$ on $V(G)\times[q]$, writing
$v_i=(v,i)$. For distinct output vertices $u_i,v_j$, set
\[
 u_iv_j\in E(S_q(G))
 \quad\Longleftrightarrow\quad
 \bigl(i=j\text{ and }u\ne v\bigr)
 \ \text{or}\ \bigl(i\ne j\text{ and }uv\in E(G)\bigr).
\]
Thus each layer $V(G)\times\{i\}$ is a clique, and the displayed rule
specifies all edges between layers. The graph has $qn$ vertices and
$q\binom n2+q(q-1)|E(G)|$ edges, so the construction has polynomial
size for fixed $q$.

Bonnet's construction~\cite[Section~3.1]{BonnetTreewidth} completes
three parts to cliques. Here the $q$ parts are copies of one source
graph. The first-branch analysis below gives a treedepth lower bound
for arbitrary $q$.

\paragraph{Notation at the first branch.}
\label{subsec:ugc-branch-notation}

Fix a nonempty graph $G$ on $n$ vertices and let $q\geq2$. Let $F$
be an elimination forest of $S_q(G)$ that is not a single chain.
Vertices in different rooted trees, or in different child subtrees
of one node, are nonadjacent. We use this property in the first-branch
analysis below.

We define its initial chain $P$ and child trees as follows. If $F$
has at least two roots, let $P$ be empty and let $F_1,\ldots,F_\nu$ be
its rooted trees. Otherwise, start at the unique root and follow the
unique child until the first vertex with at least two children.
Such a vertex exists because $F$ is not a single chain. Let $P$
include this vertex, and let $F_1,\ldots,F_\nu$ be the trees rooted at
its children. In both cases $\nu\geq2$. Measure the
depths in each $F_j$ from its own root at depth one.

For each $j\in[\nu]$, put $N_j=|V(F_j)|$ and define
\[
\begin{aligned}
 I_j&=\{i\in[q]:F_j\text{ contains a vertex from layer }i\},
 &q_j&=|I_j|,\\
 U_j&=\{v\in V(G):F_j\text{ contains a copy of }v\},
 &n_j&=|U_j|.
\end{aligned}
\]
Thus $U_j$ is the set of source vertices represented in child
$F_j$, without counting repeated copies. Of the $q_jn_j$ pairs in
$U_j\times I_j$, the child contains $N_j$ pairs. Write
$M_j=q_jn_j-N_j$ for the number of missing pairs, and let $D$
be the set of source vertices
belonging to at least two of the sets $U_j$.
We use this notation throughout the structural treedepth proof.

\subsubsection{Treedepth bounds}

\begin{theorem}\label{thm:ugc-hardness}
For every fixed real $1\leq r<4/3$, approximating treedepth within
factor $r$ on simple unweighted graphs is Unique Games-hard.
\end{theorem}

\begin{lemma}[YES upper bound]\label{lem:ugc-completeness}
For every graph $G$ on $n$ vertices and integer $q\geq1$,
\[
  \td(S_q(G))\leq q\bigl(n-\alphaI(G)\bigr)+\alphaI(G)
  =qn-(q-1)\alphaI(G).
\]
\end{lemma}

\begin{proof}
Fix a maximum independent set $I$ of $G$. Put the $q$ copies of every
vertex of $V(G)\setminus I$ on one chain. Beneath this chain, put one
chain containing the copies of $I$ in each layer. There are no edges
between the latter chains: within a layer all vertices are comparable,
and across different layers an edge would require two vertices of $I$
to be adjacent in $G$. Every edge incident to a vertex on the top chain
has comparable endpoints. This is an elimination forest of height
$q(n-|I|)+|I|$. If the top chain is empty, use the $q$ lower chains as
separate rooted trees.
\end{proof}

Write $A_{q,\eta}:=(q+1)/2+(q-1)\eta$. In a YES instance of
Theorem~\ref{thm:kr}, Lemma~\ref{lem:ugc-completeness} gives
\begin{equation}\label{eq:ugc-yes}
 \td(S_q(G))\leq qn-(q-1)\alphaI(G)
 \leq A_{q,\eta}n.
\end{equation}

\paragraph{The structural lower bound.}

The clique layers force all surviving vertices of a layer into a
single child subtree at every branch. Moreover, source vertices that
occur in more than one child form an independent set. These two
observations allow induction on the number of layers. The following
inequality combines the resulting child bounds without assuming that
the children have equal sizes.

\begin{lemma}[Branching inequality]\label{lem:branching-ineq}
For a positive integer $\nu$, positive reals $q_1,\ldots,q_\nu$, and nonnegative reals
$n_1,\ldots,n_\nu$,
\[
  3\sum_{j=1}^\nu q_jn_j-2\max_j(q_jn_j)
  \leq \left(\sum_{j=1}^\nu q_j\right)\left(\sum_{j=1}^\nu n_j\right).
\]
\end{lemma}

The proof is given in Appendix~\ref{app:branching-proof}.

\begin{lemma}[Layers in different children]\label{lem:ugc-layer-usage}
The sets $I_j$ are pairwise disjoint. Moreover,
$\nu\leq\sum_jq_j\leq q$ and $1\leq q_j\leq q-\nu+1<q$ for every
$j\in[\nu]$.
\end{lemma}

\begin{proof}
Vertices in different child trees are incomparable, so a clique
layer cannot occur in two children. Thus the $I_j$ are pairwise
disjoint. Each is nonempty, giving $\nu\leq\sum_jq_j\leq q$.
The other $\nu-1$ children use at least $\nu-1$ layers, so
$q_j\leq q-\nu+1<q$.
\end{proof}

\begin{lemma}[Source vertices in different children]
\label{lem:ugc-source-supports}
The set $D$ is independent in $G$, and
\begin{equation}\label{eq:ugc-support-count}
 \sum_jn_j\leq n+(\nu-1)\alphaI(G)
 \leq n+(q-1)\alphaI(G)\leq qn.
\end{equation}
\end{lemma}

\begin{proof}
Take distinct $u,v\in D$. Choose a copy of $u$ in one child. Since
$v$ occurs in at least two children, some copy of $v$ occurs in a
different child. The two copies have different layer indices by
Lemma~\ref{lem:ugc-layer-usage}. If $uv$ were an edge of $G$, the
definition of $S_q(G)$ would make these copies adjacent. This
contradicts their incomparability. Hence $uv\notin E(G)$, proving
that $D$ is independent and $|D|\leq\alphaI(G)$.

The first occurrence of every source vertex contributes at most one
to $\sum_jn_j$. Only vertices in $D$ contribute further occurrences,
and each contributes at most $\nu-1$ of them. Thus,
\[
 \sum_jn_j
 =\Bigl|\bigcup_jU_j\Bigr|
   +\sum_{v\in D}\bigl(|\{j:v\in U_j\}|-1\bigr)
 \leq n+(\nu-1)|D|\leq n+(\nu-1)\alphaI(G).
\]
The remaining inequalities in~\eqref{eq:ugc-support-count} follow
from $\nu\leq q$ and $\alphaI(G)\leq n$.
\end{proof}

\begin{lemma}[The full graph associated with a child]
\label{lem:ugc-child-restriction}
After relabeling the layers in $I_j$ by $[q_j]$, the graph induced by
$V(F_j)$ is obtained from $S_{q_j}(G[U_j])$ by deleting $M_j$
vertices. In particular, $M_j\geq0$.
\end{lemma}

\begin{proof}
All vertices of $F_j$ lie in $U_j\times I_j$. On this set, the
adjacency rule for $S_q(G)$ agrees with that of $S_{q_j}(G[U_j])$
after relabeling the layers. The child contains $N_j$ of the
$q_jn_j$ possible vertices, leaving $M_j=q_jn_j-N_j\geq0$ missing.
\end{proof}

\begin{lemma}[Prefix and child heights]
\label{lem:ugc-forest-accounting}
The prefix and the children satisfy
$|P|+\sum_jN_j=qn$ and $h(F)=|P|+\max_jh(F_j)$.
\end{lemma}

\begin{proof}
By the construction of $P$, its vertex set and the sets $V(F_j)$
partition $V(F)$, proving the first identity. A longest path in $F$
consists of $P$ followed by a longest path in one child, proving
the second.
\end{proof}

\begin{lemma}[Restoring deleted vertices: height]\label{lem:restore-vertices}
Let $H$ be a graph on $N$ vertices, let $X\subseteq V(H)$ have size
$d$, and let $F$ be an elimination forest of $H-X$. There is an
elimination forest $F'$ of $H$ with $h(F')=h(F)+d$.
In particular, $\td(H)\leq\td(H-X)+d$.
\end{lemma}

\begin{proof}
Place $X$ on a chain above all roots of $F$. If $X$ is empty,
keep $F$. Every new vertex is comparable with every vertex. Hence
$H$ has an elimination forest whose height increases
by $d$. Minimizing height gives the inequality.
\end{proof}

\begin{theorem}[Structural treedepth lower bound]\label{thm:ugc-structural}
For every finite simple graph $G$ on $n$ vertices and every integer
$q\geq1$,
\begin{equation}\label{eq:ugc-structural}
  \td(S_q(G))\geq
  \frac{2q}{3}n-\frac{q(q-1)}{3}\alphaI(G).
\end{equation}
\end{theorem}

\begin{proof}
We prove the statement by induction on $q$, with the induction
hypothesis asserted for every finite simple source graph.

\emph{Base cases and the first branch.}
If $G$ is empty, then $S_q(G)$ is empty and
both sides of~\eqref{eq:ugc-structural} are zero. For $q=1$,
$S_1(G)$ is a clique, so $\td(S_1(G))=n\geq2n/3$.

Now fix $q\geq2$, assume the statement for every smaller number of
layers, and take an arbitrary elimination forest $F$ of $S_q(G)$. We
will lower-bound its height $h(F)$. A single chain has height
$qn\geq2qn/3-q(q-1)\alphaI(G)/3$, so assume that $F$ is not a
single chain.

Use the first-branch notation from
Section~\ref{subsec:ugc-branch-notation}. By
Lemmas~\ref{lem:ugc-layer-usage}, \ref{lem:ugc-source-supports},
and~\ref{lem:ugc-forest-accounting}, we have
$|P|=qn-\sum_jN_j$, $1\leq q_j<q$, and
$\sum_jn_j\leq n+(\nu-1)\alphaI(G)$.
The strict inequality $q_j<q$ is what permits induction within every
child.

\emph{Induction within each child.}
By Lemma~\ref{lem:ugc-child-restriction}, the graph induced by
child $j$ is $S_{q_j}(G[U_j])$ with $M_j$ vertices deleted. Let $X_j$
be this deleted set. Thus $F_j$ is an elimination forest of
$S_{q_j}(G[U_j])-X_j$. Lemma~\ref{lem:restore-vertices} and the height
of $F_j$ give
\[
 \td(S_{q_j}(G[U_j]))
 \leq \td(S_{q_j}(G[U_j])-X_j)+M_j
 \leq h(F_j)+M_j.
\]
Induction and $\alphaI(G[U_j])\leq\alphaI(G)$ therefore give
\[
  h(F_j)\geq\frac{2}{3}q_jn_j-\frac{q_j(q_j-1)}3\alphaI(G)-M_j.
\]

\emph{Combining the child bounds.}
Choose an index $k$ maximizing $q_jn_j$. By
Lemma~\ref{lem:ugc-layer-usage}, $q_k\leq q-\nu+1$. Consequently,
$q_k(q_k-1)\leq(q-\nu+1)(q-\nu)$, so
$h(F_k)\geq\frac23\max_j(q_jn_j)
-\frac{(q-\nu+1)(q-\nu)}3\alphaI(G)-M_k$.
The height identity in Lemma~\ref{lem:ugc-forest-accounting} gives
$h(F)\geq |P|+h(F_k)$. Using $|P|=qn-\sum_jN_j$ and
$N_j=q_jn_j-M_j$ leaves the missing-vertex contribution
$\sum_jM_j-M_k\geq0$. Apply Lemma~\ref{lem:branching-ineq} to
the resulting first-branch bound. The last line also uses
$\sum_jq_j\leq q$ and~\eqref{eq:ugc-support-count}:
\[
\begin{aligned}
 h(F)&\geq qn-\sum_jq_jn_j+\frac23\max_j(q_jn_j)
       -\frac{(q-\nu+1)(q-\nu)}3\alphaI(G)\\
 &=qn-\frac13\left(3\sum_jq_jn_j-2\max_j(q_jn_j)\right)
       -\frac{(q-\nu+1)(q-\nu)}3\alphaI(G)\\
 &\geq qn-\frac13\left(\sum_jq_j\right)\left(\sum_jn_j\right)
       -\frac{(q-\nu+1)(q-\nu)}3\alphaI(G)\\
 &\geq\frac{2q}{3}n-
    \frac{q(\nu-1)+(q-\nu+1)(q-\nu)}3\alphaI(G).
\end{aligned}
\]
The error coefficient satisfies
\[
 q(\nu-1)+(q-\nu+1)(q-\nu)
 =q(q-1)-(\nu-1)(q-\nu)\leq q(q-1),
\]
because $2\leq \nu\leq q$. Substitution proves
\eqref{eq:ugc-structural} for every $F$, completing the induction.
\end{proof}

For $q\geq2$, write $B_{q,\eta}:=2q/3-q(q-1)\eta/3$.
For a NO instance of Theorem~\ref{thm:kr}, the structural bound gives
\begin{equation}\label{eq:ugc-no}
 \td(S_q(G))\geq\frac{2q}{3}n-\frac{q(q-1)}3\alphaI(G)
 >B_{q,\eta}n.
\end{equation}

\subsubsection{The approximation gap}

\begin{proof}[Proof of Theorem~\ref{thm:ugc-hardness}]
Fix $1\leq r<4/3$ and choose an integer $q>3r/(4-3r)$.
For all sufficiently small rational $\eta\in(0,1/10)$,
\[
\begin{aligned}
 B_{q,\eta}-rA_{q,\eta}
 &=\left(\frac{2q}{3}-\frac{q(q-1)}3\eta\right)
   -r\left(\frac{q+1}{2}+(q-1)\eta\right)\\
 &=\frac{q(4-3r)-3r}{6}
   -(q-1)\left(\frac q3+r\right)\eta>0.
\end{aligned}
\]
The choice of $q$ makes the constant term positive, and then a
sufficiently small fixed rational $\eta$ keeps the difference positive.
Fix such an $\eta$. Equations~\eqref{eq:ugc-yes} and~\eqref{eq:ugc-no}
give a gap greater than $r$. Since $q$ and $\eta$ depend only
on $r$, this is a polynomial reduction.
\end{proof}
